\section{The Big Picture}

\subsection{The problem}

A university lab has dozens of Windows desktops with NVIDIA graphics cards.
Most of the day they sit idle. Meanwhile students training neural networks for
coursework or a thesis have no GPU at all. distributetr9 connects the two:

\begin{enumerate}
\item A student uploads \texttt{train.py} (plus optional \texttt{requirements.txt})
      in a web portal and asks for, say, ``a GPU with at least 6 GB''.
\item A central server, the \textbf{control plane}, stores the job and puts it
      in a queue.
\item A small program on every lab PC, the \textbf{agent}, keeps asking the
      server ``do you have work for me?''. A free PC with a matching GPU gets
      the job.
\item The agent runs the script inside a locked-down \textbf{Docker
      container}, streams the output lines back live, and uploads the files
      the script wrote.
\item The student watches the logs in the browser and downloads the results.
\end{enumerate}

That sounds simple. The hard part, and the reason most of the code exists, is
what happens when things go wrong: a PC reboots in the middle of a job, Wi-Fi
drops for a minute, a student cancels at the exact moment the job finishes, or
two PCs both think they own the same job. The project's specification puts it
in one sentence: \emph{the hard problem is not the portal, it is keeping exactly
one accepted result when nodes disconnect, leases expire, completions are
retried, and cancellation races completion.}

\subsection{The components}

\begin{verbatim}
   Student's browser
        |  HTTPS
        v
  +-----------+      +---------------------------------------------+
  |  Portal   | ---> |              Control plane (FastAPI)        |
  | (Next.js) |      |  REST API  |  scheduler  |  lease worker    |
  +-----------+      +------+-------------+-----------------+------+
                            |             |                 |
                      +-----v----+  +-----v-----+    +------v------+
                      | Postgres |  |   MinIO   |    |  Registry   |
                      | (state)  |  | (files,   |    | (container  |
                      |          |  |  logs)    |    |  images)    |
                      +----------+  +-----------+    +------^------+
                            ^                               |
              HTTP: register, heartbeat, claim,             | docker pull
              renew, logs, outputs, complete                |
                            |                               |
   +------------------------+----------+--------------------+------+
   |  Lab PC 1 (Windows + WSL Ubuntu)  |  Lab PC 2  ...  Lab PC 60 |
   |  launcher -> agent -> docker run  |                           |
   +-----------------------------------+---------------------------+
\end{verbatim}

\begin{description}
\item[Portal] A Next.js web app. It never talks to the database; it forwards
  every request to the control plane's API.
\item[Control plane] A Python FastAPI service. It owns all state and makes
  every decision: who may log in, which job goes to which PC, whether a result
  is accepted.
\item[PostgreSQL] The source of truth for users, jobs, attempts and nodes.
  Its row locks are what make the scheduler safe when many PCs ask for work at
  once.
\item[MinIO] An S3-compatible object store for submitted archives, log
  chunks and output files.
\item[Registry] A Docker image registry (\texttt{registry:2}). Custom Python
  environments are built once and pulled from here by every PC.
\item[Agent] A Python program with no third-party dependencies that runs in
  WSL Ubuntu on each lab PC. It registers, heartbeats, claims work and runs it.
\item[Launcher and supervisor] A Windows scheduled task starts WSL at boot
  without anyone logging in; a small bash launcher keeps the agent running and
  rolls back bad updates.
\end{description}

\subsection{Two kinds of node: trainers and builders}

Every lab PC runs the same agent program, in one of two roles chosen at install
time:

\begin{description}
\item[Training nodes] (\texttt{role: train}, the default) run students' jobs.
  They never build environments.
\item[Build nodes] (\texttt{role: build}) build custom environments from
  \texttt{requirements.txt}, test them and push them to the registry. They never
  run training jobs.
\end{description}

The agent reports its role when it registers, and the same claim endpoint
hands out different work depending on it:

\begin{lstlisting}[language=Python]
# control/app/routers/agent.py (simplified)
@router.post("/agent/claim")
async def claim_job(node: Node = Depends(get_current_node), ...):
    if is_train_node(node):
        attempt = await claim(session, node)          # a training job
        ...
    else:
        build = await claim_build(session, node)      # an environment build
        ...
\end{lstlisting}

Why separate them? A build downloads packages from the internet, compiles
wheels and pushes multi-gigabyte images; a training job wants a quiet machine
with its whole GPU. Keeping them apart means a slow \texttt{pip install} never
delays someone's training run, and only the builder needs internet access and
registry push credentials. In this repository the role is chosen with
\verb|--role build| on the agent or installer, \verb|-Role build| on the
Windows bootstrap, or \verb|LAB_NODE_ROLE=build| in the environment.

\paragraph{A current limitation.} A builder only takes builds for its own GPU
class, or builds marked \texttt{cpu}. Jobs that ask for a GPU class and bring a
\texttt{requirements.txt} produce a build for that class, so a CPU-only builder
cannot serve them; they wait in \texttt{building} until a builder with that GPU
class exists. Catalog environments need no builder at all.

\subsection{Life of a job}

Following one job end to end is the fastest way to understand the system.

\begin{enumerate}
\item \textbf{Submit.} The portal sends \texttt{POST /jobs} with the archive.
  The control plane checks the archive, stores it in MinIO under
  \texttt{submissions/<job-id>/input.tar.gz}, and inserts a \texttt{jobs} row
  with \texttt{desired\_state = run}.
\item \textbf{Claim.} Every agent calls \texttt{POST /agent/claim}. The call is
  a \emph{long poll}: the server holds it open for up to 25 seconds until a job
  fits. When one fits, the server inserts a \texttt{job\_attempts} row with a
  fresh random \textbf{fencing token} and a \textbf{lease} that expires in 30
  seconds.
\item \textbf{Run.} The agent downloads the input, starts \texttt{docker run},
  and every few seconds calls \texttt{/renew} to extend the lease. Log lines are
  posted in numbered chunks.
\item \textbf{Upload.} Output files are uploaded under an
  \emph{attempt-specific} prefix, not the final location.
\item \textbf{Complete.} The agent calls \texttt{/complete} with its fencing
  token. The server checks the token is still the current one, verifies every
  output checksum, copies the files to \texttt{jobs/<id>/promoted/}, and marks
  this attempt as \emph{the} accepted result.
\item \textbf{Failure path.} If the PC vanishes, the lease runs out. A
  background worker marks the attempt \texttt{unknown}, then \texttt{failed}
  with category \texttt{lease\_lost}, and the job goes back in the queue for
  another PC. If the old PC comes back and tries to complete, its token no
  longer matches and the server answers \texttt{409 fenced}.
\end{enumerate}

\subsection{Words used throughout}

\begin{description}
\item[Node] A lab PC running the agent.
\item[Job] What the student asked for. A job can have several attempts.
\item[Attempt] One try at running a job on one node. At most one attempt per
  job is ever \emph{accepted}.
\item[Lease] A time-limited claim on a job. Renew it or lose it.
\item[Fencing token] A random value that identifies the current attempt; stale
  attempts are rejected because they hold an old token.
\item[Compatibility class] A label like \texttt{ampere-8g} or
  \texttt{turing-6g} describing a GPU's architecture and memory band.
\item[Environment] A container image a job runs in: a maintained catalog image
  or one built from the student's \texttt{requirements.txt}.
\end{description}
